\makeabstract
    {
     Investigating Liquid-Liquid Phase Separation between Nucleoprotein and dsDNA 
    }
    {
    Knowledge Marikiti\textsuperscript{1}, Linxi Cai\textsuperscript{1}, Ashley Carter\textsuperscript{1}
    }
    {
    \textsuperscript{1} Biochemistry \& Biophysics Program, Amherst College, Amherst, MA
    }
    {
    Vaccines face two main limitations: they rely on a person{\textquoteright}s immune system and on a low mutation rate. Broad-spectrum antivirals avoid this by targeting conserved viral pathways, such as genome packaging. In genome packaging, nucleocapsid (N) proteins and RNA form membrane-less compartments using liquid-liquid phase separation (LLPS). This phase-separated compartment increases the local concentration of N protein and RNA, causing mass production of nucleocapsids. Our goal is to determine how this LLPS compartment forms. One hypothesis is that formation occurs through non-specific interactions between the positively charged N-protein and the negatively charged RNA backbone. Understanding this underlying physics may provide insight into how to design effective antivirals.

    }
    {
    To test whether the interactions are non-specific, we pipetted a solution containing N-protein and double-stranded DNA, rather than RNA, into a sample chamber and used microscopy to visualize the formation of LLPS compartments.

    }
    {
    To test whether stoichiometry could drive condensation, we evaluated N-protein and dsDNA mixtures at a constant 2:1 molar ratio across multiple concentration tiers (90:45 nM, 70:35 nM, and 60:30 nM). However, no liquid-liquid phase separation (LLPS) was observed under any of these tested conditions 
    }
    {
     All the solutions were mixed with phosphate buffer to prevent unwanted N-protein aggregation. The theoretical 2:1 sweet spot did not produce significant results. Some aggregates were observed under brightfield microscopy but did not fluoresce under fluorescence microscopy, suggesting they are not LLPS. The lack of observable phase separation with dsDNA in this experiment highlights the potential existence of hidden energetic or kinetic barriers. Rather than disproving LLPS capability, these findings suggest that dsDNA phase behavior is governed by multi-faceted interactions beyond basic electrostatic affinity, requiring further exploration of higher-order structural or thermodynamic determinants. 

    }

    \newpage
    